\section{Limitations}

The external cohort is modest relative to the number of sources of EEG
heterogeneity. HBN and MIPDB differ in recruitment, acquisition hardware,
recording procedures, demographic composition, and potentially data quality.
External performance therefore combines representation accessibility with
domain shift; it cannot isolate a single causal source of error.

The inference covers ten optimization seeds and four prespecified heads. It
does not cover all initializations, regularization strengths, checkpoint rules,
architectures, REVE layers, or pooling functions. The bootstrap quantifies
uncertainty under the implemented seed-and-subject resampling scheme, while the
randomization test treats seed as the sign-flip unit. With only ten seeds, the
randomization distribution is necessarily discrete.

The MIPDB primary analysis is restricted to ages within HBN training support.
The extrapolation cohort was not included in the confirmatory comparison, so no
claim is made about out-of-support ages. The published REVE source list does not
name MIPDB among pretraining datasets, but absence from that list is evidence,
not a cryptographic guarantee of encoder-level independence.

Chronological-age decoding should not be confused with validation of a clinical
brain-age biomarker. Negative $R^2$ and non-ideal calibration in this external
cohort show that correlation can coexist with substantial absolute error and
systematic scaling. Finally, the present result does not establish equivalence,
a true zero effect, or the absence of benefit from an untested method.
